\mxunit{premiere-c07}{Géométrie repérée : droites, projections et cercles}\label{C07}
\begin{keypoint}[Le parcours du chapitre]
\textbf{Objectifs.} Utiliser un vecteur normal ; calculer un projeté ; déterminer ou reconnaître un cercle ; traduire une configuration en équations.

\textbf{Prérequis.} Coordonnées vectorielles, produit scalaire, systèmes linéaires, complétion du carré. Repère toujours orthonormé.

\textbf{Repère annuel.} 5, 7 et 12 janvier : 6 h\end{keypoint}
\section*{Diagnostic : \mxcount{diag} questions, 10 minutes}\label{C07-D}
Répondre sans consulter le cours. Indiquer les questions pour lesquelles une aide est nécessaire. Ce diagnostic ne donne pas lieu à une note.
\begin{enumerate}
\item $\pt{2;-1}\cdot\pt{1;2}$ ?
\item Le point $\pt{1;2}$ vérifie-t-il $2x-y=0$ ?
\item Résoudre $x+y=5$, $x-y=1$.
\item Distance de $\pt{1;2}$ à $\pt{4;6}$ ?
\item Compléter $x^2-4x$.
\item Résoudre $x^2+y^2=0$ dans $\R^2$.
\mxcounted{diag}
\end{enumerate}
\paragraph{Après le diagnostic.} Reprendre les notions signalées avec le professeur ; refaire une question analogue, puis expliquer la règle utilisée.
\paragraph{Oral et logique.} Présenter les deux conditions caractérisant un projeté. Donner une preuve géométrique et une preuve algébrique d’une médiatrice.
\mxnewpage
\section*{Chercher avant de formaliser}\label{C07-RECH}
À partir de $P\pt{4;3}$ et $d:y=2x$, comparer les distances de $P$ à trois points choisis sur $d$. Chercher ensuite le point de $d$ le plus proche de $P$ parmi tous les points de la droite, puis justifier que la distance obtenue est minimale.
\paragraph{Organisation.} Recherche individuelle 5 min, confrontation à deux 5 min, mise en commun 5 min. Garder une trace d’un essai, même non abouti. Une conjecture doit ensuite être justifiée.
\mxnewpage
\section{Vecteur normal et équation cartésienne}\label{C07-U1}
\begin{keypoint}[Définition et résultats]
Un vecteur normal non nul à une droite est orthogonal à tout vecteur directeur de cette droite. Si $\vec n=\pt{a;b}\ne\pt{0;0}$ est normal et $A\pt{x_A;y_A}$ appartient à d, alors $d$ a pour équation $a(x-x_A)+b(y-y_A)=0$, soit $ax+by+c=0$. Un vecteur directeur possible est $\pt{-b;a}$.
\end{keypoint}
\paragraph{Statut et justification.} $M\in d$ équivaut à $\vect{AM}\cdot\vec n=0$. Dans un repère orthonormé, développer ce produit donne l’équation. Réciproquement l’ensemble ainsi défini est la droite passant par A orthogonale à n.
\begin{example}Par $A\pt{1;2}$ et normale $\pt{2;-3}$ : $2(x-1)-3(y-2)=0$, donc $2x-3y+4=0$. Le vecteur $\pt{3;2}$ dirige la droite : produit $6-6=0$.\end{example}\label{C07-EX1}
\begin{pitfall}$\pt{a;b}$ est normal, et non directeur en général. $\pt{a;b}$ ne doit pas être nul.\end{pitfall}
\begin{application}Équation de la droite par $\pt{2;-1}$ de normale $\pt{3;4}$.\end{application}\label{C07-AP1}
\section{Projeté orthogonal et distance}\label{C07-U2}
\begin{keypoint}[Définition et résultats]
Le projeté H de P sur d est caractérisé par $H\in d$ et $\vect{PH}$ colinéaire à un vecteur normal de d. Pour $d:ax+by+c=0$, écrire $H=P-\lambda\pt{a;b}$, avec $\lambda =\dfrac{ax_P+by_P+c}{a^2+b^2}$. La distance est $PH=\dfrac{|ax_P+by_P+c|}{\sqrt {a^2+b^2}}$. La formule de distance est un complément déduit, pas un automatisme isolé.
\end{keypoint}
\paragraph{Statut et justification.} Substituer $x_H=x_P-\lambda a$, $y_H=y_P-\lambda b$ dans l’équation : $ax_P+by_P+c-\lambda(a^2+b^2)=0$. Le dénominateur est positif ; on obtient H puis $PH=|\lambda|\sqrt{a^2+b^2}$.
\begin{example}Pour $P\pt{3;4}$ et $d:x+y-2=0$, $\lambda =\dfrac{3+4-2}{2}=\dfrac{5}{2}$, donc $H\pt{\dfrac{1}{2};\dfrac{3}{2}}$ et $PH=\dfrac{5}{\sqrt 2}=\dfrac{5\sqrt 2}{2}$. Vérifier $x_H+y_H=2$.\end{example}\label{C07-EX2}
\mxinclude{figures/C07-U2}
\begin{pitfall}La distance est positive ou nulle : conserver la valeur absolue. Deux vérifications, appartenance et perpendicularité, sécurisent le projeté.\end{pitfall}
\begin{application}Projeter $P\pt{4;1}$ sur $d:x=2$.\end{application}\label{C07-AP2}
\section{Équation d’un cercle}\label{C07-U3}
\begin{keypoint}[Définition et résultats]
Le cercle de centre $I\pt{\alpha;\beta}$ et de rayon $r>0$ a pour équation $(x-\alpha)^2+(y-\beta)^2=r^2$. Une équation $x^2+y^2+ux+vy+w=0$ se traite par complétion des deux carrés. Si le rayon au carré obtenu est négatif, l’ensemble est vide ; s’il est nul, c’est un point.
\end{keypoint}
\paragraph{Statut et justification.} $M$ appartient au cercle si et seulement si $IM=r$. Les deux membres étant positifs ou nuls, on peut élever au carré et utiliser la formule de distance en repère orthonormé.
\begin{example}$x^2+y^2-4x+6y+9=0$ devient $(x-2)^2+(y+3)^2=4$. Centre $\pt{2;-3}$, rayon 2. Les signes des coordonnées du centre sont opposés à ceux figurant dans les parenthèses.\end{example}\label{C07-EX3}
\begin{pitfall}Une équation quadratique en x et y n’est pas toujours un cercle : vérifier les coefficients et le rayon réel.\end{pitfall}
\begin{application}Reconnaître $x^2+y^2+2x-8y+17=0$.\end{application}\label{C07-AP3}
\section{Résoudre une configuration}\label{C07-U4}
\begin{keypoint}[Définition et résultats]
Une intersection est obtenue en résolvant simultanément les équations. Pour un cercle et une droite parallèle à un axe, substituer la coordonnée constante donne une équation du second degré. Pour une médiatrice, écrire l’égalité des distances au carré ; les termes quadratiques s’éliminent.
\end{keypoint}
\paragraph{Statut et justification.} Les équivalences reposent sur l’appartenance aux deux ensembles. L’égalité $MA=MB$ est équivalente à $MA^2=MB^2$, puisque les distances sont positives.
\begin{example}Le cercle $x^2+y^2=25$ et la droite $x=3$ se coupent en $\pt{3;4}$ et $\pt{3;-4}$. La médiatrice de $A\pt{0;0}$ et $B\pt{4;2}$ vérifie $x^2+y^2=(x-4)^2+(y-2)^2$, donc $2x+y=5$.\end{example}\label{C07-EX4}
\begin{pitfall}Une intersection n’est pas l’union des ensembles de solutions. Vérifier les points dans les deux équations.\end{pitfall}
\begin{application}Intersecter $(x-1)^2+(y+2)^2=9$ avec $y=1$.\end{application}\label{C07-AP4}
\mxnewpage
\section*{Parcours d’exercices}\label{C07-PARC}
Commencer par les exercices 1 et 2. Choisir ensuite un exercice 3 ou 4 sans aide. Les exercices 5 à 7 sont une réserve de transfert, de recherche ou de programmation : le professeur en désigne un. L’exercice 8 se reprend à distance.
\begingroup
\mxexercisehead{C07-E1}{Construire une droite}
Le plan est muni d’un repère orthonormé. La droite $d$ passe par $A\pt{-1;3}$ et admet le vecteur $\vec n=\pt{2;5}$ pour vecteur normal.
\begin{enumerate}[label=\arabic*.]
\item Écrire la condition d’appartenance d’un point $M\pt{x;y}$ à $d$ à l’aide d’un produit scalaire. En déduire une équation cartésienne de $d$.
\item Donner un vecteur directeur non nul de $d$, puis les coordonnées d’un point de $d$ distinct de $A$.
\item Vérifier votre point par substitution dans l’équation.
\end{enumerate}
\mxexercisehead{C07-E2}{Calculer un projeté}
Dans un repère orthonormé, on donne le point $P\pt{4;3}$ et la droite $d$ d’équation $2x-y=0$. On note $H$ le projeté orthogonal de $P$ sur $d$.
\begin{enumerate}[label=\arabic*.]
\item Déterminer les coordonnées de $H$ en traduisant les deux conditions : $H$ appartient à $d$ et $\vect{PH}$ est normal à $d$.
\item Vérifier séparément ces deux conditions avec les coordonnées obtenues.
\item Calculer la distance de $P$ à la droite $d$ sous forme exacte.
\end{enumerate}
\mxexercisehead{C07-E3}{Reconnaître un cercle}
On considère l’ensemble des points $M\pt{x;y}$ du plan orthonormé dont les coordonnées vérifient\[x^2+y^2-6x+4y-12=0.\]
\begin{enumerate}[label=\arabic*.]
\item Montrer que cet ensemble est un cercle et préciser son centre et son rayon.
\item Déterminer tous ses points d’intersection avec la droite d’équation $y=-2$.
\item Interpréter géométriquement la position de ces deux points sur le cercle.
\end{enumerate}
\mxexercisehead{C07-E4}{Médiatrice}
Dans un repère orthonormé, on donne $A\pt{-1;1}$ et $B\pt{3;3}$. On cherche une équation de la médiatrice de $[AB]$.
\begin{enumerate}[label=\arabic*.]
\item Déterminer cette équation en utilisant les propriétés géométriques de la médiatrice. Expliquer le choix d’un point et d’un vecteur normal.
\item Vérifier que le milieu de $[AB]$ satisfait l’équation obtenue.
\end{enumerate}
\mxexercisehead{C07-E5}{Le meilleur accès}
Une route rectiligne est représentée par la droite $d$ d’équation $y=2x$ dans un repère orthonormé dont l’unité de longueur vaut 1 km. Un bâtiment se situe au point $P\pt{4;3}$. On veut construire un accès rectiligne reliant le bâtiment à un point de la route.
\begin{enumerate}[label=\arabic*.]
\item Déterminer le point de la route qui rend cet accès le plus court possible. Calculer la longueur exacte de l’accès.
\item Justifier que votre choix est optimal parmi tous les points de la route, en comparant les distances à l’aide d’un triangle rectangle. Une seule observation sur un dessin ne suffit pas.
\end{enumerate}
\mxexercisehead{C07-E6}{Paramètre et ensemble}
Pour un paramètre réel $k$, on définit $\mathcal E_k$ comme l’ensemble des points $M\pt{x;y}$ dont les coordonnées vérifient\[x^2+y^2-2x+4y+k=0.\]
\begin{enumerate}[label=\arabic*.]
\item Déterminer la nature de $\mathcal E_k$ selon la valeur de $k$. Organiser la réponse en cas disjoints.
\item Lorsque cet ensemble est un cercle, préciser son centre et son rayon ; dans les autres cas, décrire exactement l’ensemble obtenu.
\end{enumerate}
\mxexercisehead{C07-E7}{Une parabole comme lieu}
\textbf{Approfondissement facultatif.} Dans un repère orthonormé, on fixe le point $F\pt{0;2}$. On cherche les points $M\pt{x;y}$ dont la distance à $F$ est égale à la distance à l’axe des abscisses.
\begin{enumerate}[label=\arabic*.]
\item Exprimer le carré de chacune des deux distances en fonction de $x$ et de $y$.
\item En déduire une équation de l’ensemble recherché. Vérifier la réciproque afin de justifier que cette équation décrit exactement tous les points cherchés.
\end{enumerate}
\mxexercisehead{C07-E8}{Reprise différée}
Dans un repère orthonormé, $\mathcal C$ est le cercle de diamètre $[AB]$, avec $A\pt{0;0}$ et $B\pt{4;0}$. On considère la droite $d$ d’équation $x=5$.
\begin{enumerate}[label=\arabic*.]
\item Déterminer si $\mathcal C$ et $d$ ont un point commun, en résolvant une équation.
\item Confirmer la réponse en comparant la distance du centre du cercle à la droite et le rayon du cercle.
\end{enumerate}
\paragraph{Travail personnel.} Effectuer uniquement la tâche désignée à l’issue de la séance, dans la progression datée : 10–20 min maximum. Les exercices dépendent des notions déjà étudiées ; certaines questions sont donc différées. Reprendre l’exercice 8 une semaine après le chapitre (10 min). Si le temps prévu est dépassé, noter la difficulté et apporter les essais ; ne pas accumuler toute la réserve.
\endgroup
\mxneed{12}\section*{À faire dans le manuel CQFD}
Ces exercices se traitent dans votre manuel. La progression indique la séance ; ils remplacent le travail personnel prévu ce jour-là (10 à 20 minutes au total).
\paragraph{CQFD, exercice 18 p. 279, questions a, b, c et d.}\label{C07-M1} Lire un vecteur normal non nul pour chaque équation, sans résoudre en y.
\paragraph{CQFD, exercice 21 p. 279, questions a et b.}\label{C07-M2} Écrire la condition d’orthogonalité à la normale, puis vérifier que le point donné appartient à la droite.
\begingroup\small
\mxnewpage\section*{QCM 1 : comprendre et choisir}\label{C07-Q1}
\mxCount{qcm1} questions pédagogiques à réponse unique. Choisir une réponse par question ; expliquer ensuite un distracteur. Ce format de chapitre ne reproduit pas le format officiel de l’EAM.
\begin{enumerate}
\item Une normale à $2x-3y+1=0$ est
\begin{enumerate}[label=\Alph*.,itemsep=2pt]
\item $\pt{3;2}$
\item $\pt{2;-3}$
\item $\pt{2;3}$
\item $\pt{1;0}$
\end{enumerate}
\item Un directeur de $x+2y=0$ est
\begin{enumerate}[label=\Alph*.,itemsep=2pt]
\item $\pt{-2;1}$
\item $\pt{1;2}$
\item $\pt{2;1}$
\item $\pt{0;2}$
\end{enumerate}
\item Centre de $(x+2)^2+(y-3)^2=4$ :
\begin{enumerate}[label=\Alph*.,itemsep=2pt]
\item $\pt{2;3}$
\item $\pt{-2;-3}$
\item $\pt{3;-2}$
\item $\pt{-2;3}$
\end{enumerate}
\item Rayon de $x^2+y^2=9$ :
\begin{enumerate}[label=\Alph*.,itemsep=2pt]
\item $9$
\item $81$
\item $3$
\item $\sqrt3$
\end{enumerate}
\item $x^2+y^2=-1$ décrit
\begin{enumerate}[label=\Alph*.,itemsep=2pt]
\item un cercle
\item l’ensemble vide
\item l’origine
\item une droite
\end{enumerate}
\mxcounted{qcm1}
\end{enumerate}
\endgroup
\begingroup\small
\mxneed{18}\section*{QCM 2 : comprendre et choisir}\label{C07-Q2}
\mxCount{qcm2} questions pédagogiques à réponse unique. Choisir une réponse par question ; expliquer ensuite un distracteur. Ce format de chapitre ne reproduit pas le format officiel de l’EAM.
\begin{enumerate}
\item Projeté de $\pt{3;4}$ sur l’axe des abscisses :
\begin{enumerate}[label=\Alph*.,itemsep=2pt]
\item $\pt{0;4}$
\item $\pt{0;0}$
\item $\pt{3;0}$
\item $\pt{4;3}$
\end{enumerate}
\item Distance de $\pt{1;5}$ à $y=2$ :
\begin{enumerate}[label=\Alph*.,itemsep=2pt]
\item $3$
\item $5$
\item $1$
\item $\sqrt{10}$
\end{enumerate}
\item Pour définir $ax+by+c=0$ comme droite, il faut
\begin{enumerate}[label=\Alph*.,itemsep=2pt]
\item $a\ne0$ seul
\item $b\ne0$ seul
\item $c\ne0$
\item $\pt{a;b}\ne\pt{0;0}$
\end{enumerate}
\item $x^2+y^2=0$ décrit
\begin{enumerate}[label=\Alph*.,itemsep=2pt]
\item un cercle de rayon 1
\item le seul point O
\item une droite
\item le plan
\end{enumerate}
\item Cercle $x^2+y^2=4$ et droite $x=2$ : nombre de points communs
\begin{enumerate}[label=\Alph*.,itemsep=2pt]
\item $0$
\item $2$
\item $1$
\item infini
\end{enumerate}
\mxcounted{qcm2}
\end{enumerate}
\subsection*{Vrai ou faux ? Justifier.}\label{C07-VF}
\begin{enumerate}
\item Le projeté orthogonal de P sur une droite est unique.
\item Toute équation $x^2+y^2+c=0$ décrit un cercle de rayon positif.
\item Une droite peut rencontrer un cercle en trois points distincts.
\end{enumerate}
\begin{keypoint}[Bilan personnel]Je donne une normale non nulle ; je vérifie le projeté ; je complète les deux carrés ; je distingue cercle, point et ensemble vide.

Pour chaque objectif, indiquer : acquis sans aide ; réussi avec aide ; à retravailler. Écrire un exemple et une erreur à éviter.\end{keypoint}
\endgroup
\ifmxteacher\mxinclude{corriges/C07}\fi
